\begin{auxlemma}[Evaluation at a frame with a short period]
\label{lem:deep-thin-evaluation}
Let \(Z\) be a component at a nonterminal lift, \(u_1\) one of its
selected quotient weights and \(y\) a frame on \(Z\) at a good
complex embedding. Write \(N_*=g_0g_1 15^e\) for the deep level,
\(\lambda=\lambda_1(y)\) for the first minimum in the polarization
metric and \(\mathcal B_{u_1,Z}(y)\) for the supremum of
\(\|g(y)\|^2/\|g\|_{2,Z}^2\) over nonzero cusp sections of weight
\(u_1\) on \(Z\), with component measure of mass at most one. Uniformly
over these components, weights and frames, for large family scale,
\[
 N_*\lambda^2<b\quad\Longrightarrow\quad
 \mathcal B_{u_1,Z}(y)\ge e^{-\phi(y)/625},\qquad
 \phi(y)=\frac{2f}{\operatorname{rd}(E)^2\lambda^2}.
\]
\end{auxlemma}
\begin{proof}
\proofstep{Choose a differential operator at the quotient weight.}
Put \(M_0=g_0g_1\), \(d_*=2fm\), \(R=\operatorname{rd}(E)\).
Use the specific column count \eqref{eq:deep-quotient-column-count} and exponents \eqref{eq:deep-quotient-exponents} for \(u_1\). The
corresponding theta weight \(G\) is obtained by the parallel shift \(w_\tau=n/2+1-u_{\tau,p}\) and central exponent \(n/2\). If an embedding
has \(q\) rows with exponents \(a_1\ge\cdots\ge a_q\), use one
\(j\)-row minor \(P_{\tau j}\) at each height. By the definition of
\(n\), all their column sets may be chosen disjoint. Put
\[
 P=\prod_{\tau}\prod_{j=1}^{q_\tau}
          P_{\tau j}^{\,a_{\tau j}-a_{\tau,j+1}},
 \qquad a_{\tau,q_\tau+1}=0 .
\]

This polynomial is harmonic since no two contracted opposite labels occur in
the same column. Put \(D=\deg P\). The positive last exponent gives
cuspidality: at a rational boundary the projected frame loses rank
in the corresponding full minor, so that its constant term vanishes,
as in Lemma~\ref{lem:marked-class-section-cuspidal}. A determinant
of size \(j\) contributes at most \(j!\le r^j\) to the sum of
absolute coefficients. Hence
\[
 1\le D\le fr(4b+1)\le b^2,\qquad \|P\|_1\le r^D .
\]

The host, level and column estimates in
\ref{cor:deep-boundary-orders-5} and
Lemma~\ref{lem:joint-theta-parameters} give, uniformly,
\[
 frn\le b,\quad R^7\le b,\quad d_*\le2b,
 \qquad b^{16}\le M_0\le b^{70}.
\]

Indeed \(f,r,n\) grow polynomially in \(x\), whereas
\(b=\exp(x^\epsilon)\); at each type prime
\((kb)^8\le g_i\), and the upper type bounds give
\(\log g_i\le35\log b\). In the joint level estimate use this
value of \(D\) and the smaller signature bound \(q_{\rm small}\le m\);
the numerical condition is
\[
 R_0^2R^{28}Dr=90(fr+D)R^{28}Dr=O(b^9)\le M_0 .
\]
The inequalities \(m^4\le f\) and \(R\ge\max(16,f^2)\) are independent of the weight.

\proofstep{Compare one actual translate with all frames on the original level.}
For the frame over \(y\) after forgetting the deep marking, use
the same lattice and form. Its first minimum is still \(\lambda\).
Choose the marked half \(\mathfrak m=(\bar\alpha)\), \(\alpha\bar\alpha=M_0\).
The isometry of Lemma~\ref{lem:scaled-light-metric} sends \(v\mapsto\bar\alpha v\) from the original lattice in the metric \(H_A\) to the half-level lattice in \(H'=H_A/M_0\). It preserves the first minimum, since
\[
 H'(\bar\alpha v,\bar\alpha v)=H_A(v,v).
\]

The isogeny pushes Haar probability forward to Haar probability, and the marked letter has squared length \(2\xi_R/M_0\) in \(H'\). Thus all \(M_0^{frn}\) translations of the
principal theta tensor are readings of cusp rules at weight \(G\).
Write \(b_\nu(z)\) for their invariant point norms multiplied by
the common Lie normalization \((2\pi\xi_R)^{-D/2}\).
Apply Lemmas~\ref{lem:light-thick-lift}, \ref{lem:light-frequency-column}
and~\ref{lem:light-block-jet-product} to this particular \(P\).
These lemmas treat arbitrary such ordered disjoint blocks. Changing to the
marked derivative costs \((2\xi_R/M_0)^{D/2}\). Parseval over the
half-level translates, followed by the maximum bound for a finite sum,
gives
\[
 \begin{aligned}
 M_0^D\sum_\nu b_\nu(y)^2
    &\ge \pi^{-D}[24(2d_*)^{d_*/2}]^{-n}
            2^{-(r+2)D}e^{-2a\phi(y)},\\
 a&=8100\pi n m^4d_*^4R^7/M_0,\\
 -\log\max_\nu b_\nu(y)&\le
    3b^3D+a\phi(y)+(D+frn)\log M_0 .
 \end{aligned}
\]

For the last inequality, \(D\ge1\), \(d_*\le2b\) and \(r,n\le b\)
bound half the logarithm of all powers of \(2,\pi,24,2d_*\) in
the first line by \(3b^3D\). Also
\[
 a\le 8100\pi\,16 b^{10}/b^{16}\le 1/2500
 \qquad(b\ge1000).
\]
This gives a nonzero rule at \(y\); fix a maximizing translate \(\nu\).

For its upper bound on the original level, let \(z_0\) be an orthonormal frame with value at least half its supremum. The rank-one
period estimate in \ref{cor:slice-short-heights-3}, applied with
spread at most \(3b^2\), gives first minimum at least
\(e^{-40L}\). Thus, with \(L=3\log b\),
\[
 2b_\nu(z_0)\ge\sup_z b_\nu(z),\qquad \lambda_1(z_0)\ge e^{-40L}.
\]

The four bounds needed here are
\[
 \log M_0\le30L,\qquad
 \log(fr)\le L,\qquad
 \log R\le3L,\qquad
 D+frn\le2b^2.
\]

Set \(h=frn\) and \(A=2f^{3/2}n^2R\).
For the standard tensor lattice a rational basis in its second
factor consists of integral vectors of lengths at most \(A\). Contraction with this basis
bounds the tensor minimum below by \(e^{-40L}/A\).
Use the submean and Cauchy inequalities in
\ref{cor:slice-short-heights-4} for the translated theta function
in dimension \(h\). A polydisc of coordinate radius
\[
 b_0=\min\{1/\sqrt{\pi h},e^{-40L}/(4A\sqrt h)\}
\]
embeds in the torus.

Its submean estimate costs
\(e^{1/2}(\sqrt\pi b_0)^{-h}\). A multi-index of degree \(D\)
costs \((\pi eD)^{D/2}\) times the coefficient sum by Cauchy's
inequality on the Gaussian function. The Lie factor cancels the
change from unit to marked coordinates. Including the factor two
for the half-supremum frame gives
\[
 \log\sup_z b_\nu(z)\le \log 2+\tfrac12
      +h\log(1/(\sqrt\pi b_0))+D\log r+\tfrac D2\log(eD)
       =O(b^2\log b).
\]
Here \(h\le b\), \(D\le b^2\), and \(\log A=O(\log b)\).
One may also retain the full finite marking count
\(M_0^{frn}\) in the supremum estimate since its logarithm has
the same bound.

Since \(e=\lfloor3\log b/\log3\rfloor\), eventually \(15^e\ge b\).
The thin-frame condition gives \(M_0\lambda^2\le1\) and hence
\[
 \phi(y)\ge b^{14},\qquad
 3b^3D+(D+frn)\log M_0+O(b^2\log b)\le \phi(y)/2500 .
\]
Together with the bound on \(a\), this proves
\(e^{-\phi(y)/1250}\sup_z b_\nu(z)\le b_\nu(y)\).

\proofstep{Return to the selected weight on the given component.}
Forget the refined marking and pull the chosen cusp rule back to
the deep level. In the adelic description the map enlarges the
level stabilizer on the right; the Hermitian lattice and the
orthonormal frame are carried by the same rational isometry.
The first minimum and every invariant point norm are preserved.
Write \(\pi\) for the forgetting map and \(e_Z\) for the idempotent of \(Z\). For the selected translate \(f_\nu\),
\[
 \|e_Z\pi^* f_\nu(z)\|=\|f_\nu(\pi z)\|\quad(z\in Z),
 \qquad (e_Z\pi^* f_\nu)|_{Z^c}=0 .
\]

The parallel determinant generator of Lemma~\ref{lem:flat-generic-generator},
raised at each embedding to power \(-w_\tau\), and the inverse central unit raised to power \(n/2\) in the weight return of
\eqref{eq:parallel-return-norm} and \ref{cor:deep-cover-on-core-10} convert the weight \(G\) to
\(u_1\). Both factors are invertible on the component and have
constant positive invariant norm there. Their product multiplies
both sides of the ratio estimate by one and the same positive
constant. The resulting section \(g\) is cuspidal and
\[
 \|g(y)\|\ge e^{-\phi(y)/1250}\|g\|_{\sup,Z}>0,\qquad
 \|g\|_{2,Z}\le\|g\|_{\sup,Z}.
\]
Squaring proves the assertion for \(\mathcal B_{u_1,Z}(y)\).
\end{proof}
